\section{Self-contained dependency certificate}
\label{app:dependencycertificate}

This appendix separates logical premises from provenance.  The conclusions of
Paper VII use only the following statements.  A reader may treat each imported
statement as an explicit assumption and verify every theorem in this paper
without consulting another Tau manuscript.

\begin{description}[style=nextline,leftmargin=2.8em,labelwidth=2.2em]
\item[\textbf{D1: body-first order.}]
There is a $C^2$ pre-body action $\mathcal A_{\rm pre}[y;B,s]$ with an occupied
stationary point $y_*$ and positive physical Hessian $H_B$.  The post-body
functional is evaluated on the frozen parameter $y_*$; post-body source
variation is not varied back into the body equation.  This is the only
body-first premise used in Eqs.~\eqref{eq:jointfunctional} and
\eqref{eq:triangular}.

\item[\textbf{D2: common typed source domain.}]
On the selected body and observer support there is one finite source domain
$\mathcal E_{*,O}$, one access map $A_O$, typed linear role maps
$R_a:\mathcal E_{*,O}\to\mathcal Y_a$, and stable resolution maps
$Q_O^{(a)}$.  No direct-sum decomposition of $\mathcal E_{*,O}$ is assumed.
Proposition~\ref{prop:moprassembly} is proved here from exactly this premise.

\item[\textbf{D3: occupied operation data.}]
An operation algebra $\mathfrak A_{\rm src}$, occupied vector $\Omega_s$ and
positive root covariance $X_0$ are given.  Whenever the full-occupation branch
is used, $\Omega_s$ is cyclic and every required $R_a$ is surjective.
Whenever the rooted-path branch is used,
$\operatorname{supp}X_0\subseteq V_s$ and every admitted $L_e$ preserves
$V_s=\operatorname{span}(\mathfrak A_{\rm src}\Omega_s)$.
Proposition~\ref{prop:occupation} proves the resulting full-occupation or
activation statement; physical realization of either premise set is not
imported.

\item[\textbf{D4: fixed quantum preparation.}]
The CTP boundary datum satisfies $\rho_{\rm in,O}\geq0$ and
$\operatorname{Tr}\rho_{\rm in,O}=1$.  It is held fixed while $g$, $A$ and
$\Phi$ are varied.  The CTP evolution and observer restriction produce a
normalized output state $\rho_O$.  CTP response is obtained by varying the
difference source before taking the physical equal-branch limit.  No Euler
variation with respect to a density matrix is used.

\item[\textbf{D5: standard Lorentzian leaf.}]
On the selected leaf the geometric and Abelian functionals satisfy the
classification hypotheses of Proposition~\ref{prop:standardforms}, with a
common unit packet.  Their Einstein--Hilbert and Maxwell shapes then follow
inside that class.  The matter functional, measure and renormalization
prescription are diffeomorphism and $U(1)$ gauge invariant and anomaly-free
for the Ward identities invoked here.  This is an assumed standard-limit
leaf, not a Tau derivation of those actions.

\item[\textbf{D6: operational effect source.}]
Each declared outcome family is a POVM, $E_a\geq0$ and $\sum_aE_a=I$, and its
source coupling is Eq.~\eqref{eq:effectcoupling}.  This premise, together with
D4, is sufficient for the normalized Born derivative in
Eq.~\eqref{eq:sourceowners1}; no stronger quantum reconstruction theorem is
used.

\item[\textbf{D7: clock-source boundary.}]
The clock character $\vartheta$ is an admitted external source coordinate of
the post-body functional.  Paper VII uses only the formal response derivative
$\delta W_O^{\rm ren}/\delta\vartheta$.  It neither imports nor proves a
duration law, global clock integrability or a Tau-specific clock distortion.

\item[\textbf{D8: once-counted calibration.}]
Every displayed sector uses one declared action unit, each occupied
contribution and counterterm has one ledger owner, and the overlap
$K_{m,m+2}$ is calculated from Eq.~\eqref{eq:overlapformula} and frozen before
the ROOT--TT comparison.  This bookkeeping premise forbids an additional
fitted gain but does not select the common unit, occupied mode pair, screen or
subtraction prescription.
